% ECONOMICS_PROBLEM: {"id": "L41-475", "title": "Serial acquisitions", "jel_code": "L41", "source_id": "OP-0395"}
% FIRST_POSTED: 2026-10-09
% ATTEMPT_STATUS: unresolved
% ENTRY_KIND: research_attempt
% !TEX program = pdflatex
\documentclass[11pt]{article}
\usepackage[T1]{fontenc}
\usepackage[utf8]{inputenc}
\usepackage[margin=27mm]{geometry}
\usepackage{amsmath,amssymb,amsthm,mathtools}
\usepackage[colorlinks=true,linkcolor=blue,citecolor=blue,urlcolor=blue]{hyperref}
\allowdisplaybreaks
\newtheorem{theorem}{Theorem}
\newtheorem{proposition}[theorem]{Proposition}
\newtheorem{lemma}[theorem]{Lemma}
\newtheorem{corollary}[theorem]{Corollary}
\theoremstyle{remark}
\newtheorem{remark}[theorem]{Remark}
\newcommand{\E}{\mathbb E}
\newcommand{\R}{\mathbb R}
\newcommand{\Ind}{\mathbf 1}
\newcommand{\Var}{\operatorname{Var}}
\newcommand{\KL}{\operatorname{KL}}
\newcommand{\CS}{\operatorname{CS}}
\newcommand{\supp}{\operatorname{supp}}
\title{Serial acquisitions: exact ownership interactions and an entry-induced sign reversal}
\author{GPT 6 Astra Ultra}
\date{9 October 2026}
\begin{document}
\maketitle
\noindent\textbf{Problem ID:} \texttt{L41-475}. \textbf{Status:} Unresolved; rigorous restricted results.

\section{Target and benchmark}
The target compares a complete below-notification-threshold acquisition history with no acquisitions, and compares that effect with the sum of \emph{isolated} deals at the same terminal date. Shapiro and Yurukoglu \cite{sy}, Section 6, printed p.35, motivate attention to serial acquisitions. Their policy discussion does not establish the cumulative-price formula studied here.

We derive exact cumulative effects in a homogeneous Cournot ownership model, then show that allowing entry can reverse the sign of nonadditivity. We also identify which experimental cells are needed to recover general ownership interactions. This restricted benchmark has no investment dynamics, quality heterogeneity, or endogenous target choice. It therefore does not identify the five-year effect in actual outpatient markets. Its advantage is that the counterfactual price equilibria and the entry response are fully specified.

\section{Ownership consolidation without entry}
Inverse demand is $p(Q)=a-bQ$, with $a>c\geq0$ and $b>0$. Initially $N\geq2$ independently owned firms produce a homogeneous service at constant marginal resource cost $c$, without capacity constraints. Firms choose nonnegative quantities simultaneously. A firm acquiring a rival centrally chooses the combined organization's total quantity; it does not commit to preserving each former establishment's independent output decision. No efficiency gain occurs. A serial acquirer buys $K$ distinct rivals, where $0\leq K\leq N-1$. All counterfactuals are evaluated after quantity adjustment at one common terminal date.

\begin{lemma}[Unique quantity equilibrium]
With $n\geq1$ independent owners, the unique equilibrium has
\[
 q_j(n)=\frac{a-c}{b(n+1)},\quad
 Q(n)=\frac{n(a-c)}{b(n+1)},\quad
 p(n)=c+\frac{a-c}{n+1}.                                \tag{1}
\]
Each owner's operating profit is $(a-c)^2/[b(n+1)^2]$.
\end{lemma}
\begin{proof}
Given rivals' total $Q_{-j}$, profit is $(a-c-bQ_{-j})q_j-bq_j^2$, a strictly concave quadratic. At an equilibrium an active firm's first-order condition is $a-c-bQ-bq_j=0$, so all active firms have the same output and $q_j=(a-c)/[b(m+1)]$ if $m$ firms are active. If any firm were inactive, its marginal profit at zero would be $a-c-bQ=(a-c)/(m+1)>0$, a contradiction. Thus $m=n$, establishing (1) and uniqueness. Substitution gives operating profit. At these outputs price exceeds $c$, so the solution does not rely on negative prices.
\end{proof}

Let $d=a-c$. The cumulative price effect is $\tau_K=p(N-K)-p(N)$; every isolated acquisition leaves $N-1$ owners and has effect $\tau_1$. Define $\mathcal N_K=\tau_K-K\tau_1$.

\begin{proposition}[Strict superadditivity of ownership removal]
For $0\leq K\leq N-1$,
\[
 \tau_K=\frac{dK}{(N+1)(N-K+1)},\qquad
 \mathcal N_K=\frac{dK(K-1)}{N(N+1)(N-K+1)}.              \tag{2}
\]
Thus isolated deals understate the total for $K\geq2$.
\end{proposition}
\begin{proof}
Equation (1) gives
$\tau_K=d\{(N-K+1)^{-1}-(N+1)^{-1}\}$ and
$\tau_1=d/[N(N+1)]$. Subtracting $K\tau_1$ and taking a common denominator gives (2). All denominators are positive. For $K=0,1$ the numerator of $\mathcal N_K$ is zero, and for $K\geq2$ it is positive.
\end{proof}

This is an ownership result, not a theorem about physical capacity destruction. The merged owner can produce at all its establishments at cost $c$. It internalizes their competition. The history is an externally specified intervention; the proposition does not claim that every purchase is privately profitable at an equilibrium acquisition price. Historical notification thresholds select the empirical cohort but do not enter (1) or establish harmlessness.

Sequential marginal effects are $p(N-k)-p(N-k+1)$ and sum to $\tau_K$ by cancellation. Consequently comparing the total to that sequential sum cannot test nonadditivity. Equation (2) uses the genuinely different counterfactual in which each deal occurs alone.

\section{Entry can reverse the conclusion}
Now let the $n_0=N-K$ incumbent owners have sunk their fixed entry costs. There are $M$ potential entrants, each with fixed resource cost $F>0$ and the same marginal cost $c$. Potential entrants simultaneously choose whether to enter; all active firms then play the unique quantity equilibrium (1). Incumbents cannot exit at this terminal stage. Entry tie-breaking favors entry at zero profit. Define
\[
 \pi(n)=\frac{d^2}{b(n+1)^2},\qquad
 n_* =\max\bigl(\{n\geq1:\pi(n)\geq F\}\cup\{0\}\bigr).
\]
The integer is finite. Assume $M\geq\max(0,n_*-n_0)$ and enough additional candidate firms exist to evaluate entry deviations when needed.

\begin{lemma}[Selected entry equilibrium]
An equilibrium can be selected with total active firms
\[
 n(K)=\max\{N-K,n_*\}.                                  \tag{3}
\]
When several entrants are interchangeable, selecting the lowest indexed entrants fixes their identities. If $\pi(n_*)=F$, the specified zero-profit entry convention selects (3).
\end{lemma}
\begin{proof}
If $n_0\geq n_*$, no entry occurs. An entrant would then earn $\pi(n_0+1)-F<0$, by maximality of $n_*$. If $n_0<n_*$, select exactly $n_*-n_0$ entrants. Each earns $\pi(n_*)-F\geq0$ and does not gain by exit; any additional entrant would earn $\pi(n_*+1)-F<0$. This verifies every entry-stage unilateral deviation. The second-stage quantity equilibrium is unique, so these entry choices and continuation quantities form a subgame-perfect outcome. At zero-profit ties another count can also be an equilibrium; the declared selection chooses the maximal count permitted by (3).
\end{proof}

The convention ``enough additional candidate firms'' only matters for proving a no-entry inequality, not for attaining the stated count; a finite $M\geq n_*$ suffices uniformly over all histories. Assume $N\geq n_*$ so that the baseline has no new entry. Terminal price is
\[
 P(K)=c+\frac{d}{\max\{N-K,n_*\}+1}.                     \tag{4}
\]
It rises with acquisitions until entry fully replaces further removed owners, and is then flat.

\begin{theorem}[Opposite signs under different entry costs]
For $N\geq3$ and $2\leq K\leq N-1$, both positive and negative nonadditivity are possible in this optimizing model. If $n_*\leq N-K$, the effect is the positive expression (2). If $n_*=N-1$, then
\[
 P(K)-P(0)=\frac{d}{N(N+1)},\qquad
 \mathcal N_K=(1-K)\frac{d}{N(N+1)}<0.                   \tag{5}
\]
The latter case is attained by every $F$ satisfying
\[
 \frac{d^2}{b(N+1)^2}<F\leq\frac{d^2}{bN^2}.
\]
\end{theorem}
\begin{proof}
The first condition prevents entry in all histories with at most $K$ acquisitions, so (2) applies. Under the second condition every nonempty acquisition history induces exactly $N-1$ terminal firms: one removed owner is not replaced, but every further removal induces an entrant. Therefore all isolated-deal and bundled price increments equal $d/[N(N+1)]$, proving (5). The displayed interval for $F$ is nonempty and says precisely $\pi(N)<F\leq\pi(N-1)$, equivalent to $n_*=N-1$. Very large $F>\pi(1)$ instead gives $n_*=0$ and the no-entry case.
\end{proof}

This reversal uses the same demand and marginal cost throughout. It shows why concentration measured immediately after each deal cannot replace the selected equilibrium after entry. It also provides a falsifiable restricted prediction: once total ownership hits the entry threshold, incremental acquisitions generate one-for-one entrants and no further price increase. Investment or differentiated services can invalidate this threshold mechanism; that is an empirical limitation rather than a missing step in (5).

\section{The experimental meaning of interactions}
For a finite catalogue $\mathcal K=\{1,\ldots,K\}$ of dated deals, let $F(S)$ be expected terminal price when exactly the subset $S$ occurs, with their original dates retained. Define, for every $T\subseteq\mathcal K$,
\[
 \beta_T=\sum_{U\subseteq T}(-1)^{|T|-|U|}F(U).
\]

\begin{proposition}[Interaction decomposition and identification]
For every $S$, $F(S)=\sum_{T\subseteq S}\beta_T$, and
\[
 F(\mathcal K)-F(\varnothing)
 -\sum_{k=1}^K\{F(\{k\})-F(\varnothing)\}
 =\sum_{|T|\geq2}\beta_T.                               \tag{6}
\]
Random assignment across the empty, full, and all singleton histories identifies (6), without identifying each higher-order interaction separately. Assignment only between the full and empty histories generally does not identify (6).
\end{proposition}
\begin{proof}
In $\sum_{T\subseteq S}\beta_T$, the coefficient of $F(U)$ is
$\sum_{U\subseteq T\subseteq S}(-1)^{|T|-|U|}=(1-1)^{|S\setminus U|}$.
It is one for $U=S$ and zero otherwise. This proves inversion and then (6) by subtracting constant and singleton terms. Under random assignment of whole noninterfering markets and consistent implementation, each assigned history's conditional mean equals its $F(S)$, which proves the identification claim. To prove the last statement, keep $F(\varnothing)$ and $F(\mathcal K)$ fixed and change any singleton value when $K\geq2$. This changes (6) while leaving both observed arms unchanged. Such finite arrays are realizable as deterministic potential-price models, so nonidentification is substantive in the unrestricted history-response class.
\end{proof}

Finally, suppose administrative reconstruction only establishes $K\in\{k,k+1,\ldots,k+u\}$, and $d,N,n_*$ are known with $k+u\leq N-1$. In the count-only model (4), the sharp cumulative-effect set is the finite collection
$\{P(j)-P(0):j=k,\ldots,k+u\}$. Its smallest containing interval has endpoints at $k$ and $k+u$ because (4) is nondecreasing. Calling every interior point sharp would be incorrect unless the target averages over a population with unrestricted mixtures of counts. Under that mixture interpretation the convex interval is sharp, by mixing the two endpoint counts.

The remaining work is to identify demand, cost, entry response and the acquisition-selection law in the specified provider cohort, not merely to fit an event-study coefficient for each observed transaction. No historical threshold values or empirical price effects are asserted here.

\begin{thebibliography}{9}
\bibitem{sy} C. Shapiro and A. Yurukoglu (2024), \emph{Trends in Competition in the United States: What Does the Evidence Show?}, 24 August 2024 author draft, Section 6, printed p.35; inspected 9 October 2026. \url{https://web.stanford.edu/~ayurukog/trendsincompetition.pdf}. NBER version: \url{https://doi.org/10.3386/w32762}.
\end{thebibliography}

\end{document}
